\section{Related Work} \label{sec:related}
% \syfC{重写内容和逻辑，补充ogb上的两个方法：LD和BoT}

\vpara{Multi-label Node Classification (MLNC).}
% raise awareness and 
To promote research on MLNC, MLGNC~\citep{zhao2023multi} provides three biological datasets and performs a detailed analysis on both datasets and existing methods.
From the methodology perspective, some methods consider using textual embeddings of labels to incorporate label proximity.
LARN~\citep{XIAO2022235} incorporates the all label embeddings during the neighbor aggregation process. Similarly, LANC~\citep{ZHOU2021115063} proposes attention mechanism to aggregates the label embeddings with each node.
Another line of work construct additional label-label network to use the label co-occurrence correlations.
ML-GCN~\citep{gao2019semisupervised} construct both node-centric and label-centric graphs to aggregate information, while MLNE~\citep{8894391} use random walk to get embedding on both label-label network and original graph.
There are also some methods that modify the approach to modeling graphs from multi-label perspective.
MLGW~\citep{8970680} leverage label-specific agents to walk on graphs. VariMul~\citep{10.1145/3459637.3482391} utilizes VGAE to derive node and label embeddings and designs a confidence ranking loss to mine pairwise label correlations.
Recent work improves MLNC by addressing the ambiguity in multi-label graph structures~\citep{bei2024correlation} and the limited expressive power for multi-label tasks~\citep{zhao2024gnn}.
% zhao2024gnn
% bei2024correlation
However, these work do not explicitly analyze and quantify the influences between labels and exploit them to enhance the MLNC task.

\vpara{Decoupled GNNs.}
% In this section, we discuss some related decoupled GNNs. 
Compared to conventional GNNs, decoupled GNNs explicitly isolate the two operations and aggregate all propagation (P) and transformation (T) operations separately for processing.
APPNP~\cite{gasteiger2018predict} and DAGNN~\cite{liu2020towards} first propose that decoupling the two operations can enable GNNs to go much deeper without causing over-smoothing. Although SGC~\cite{wu2019simplifying} and lightGCN~\cite{he2020lightgcn} do not discuss the decouple paradigm, it demonstrate superior efficiency and performance. PTA~\cite{dong2021equivalence} further explore the paradigm from the perspective of label propagation. 
Since the APPNP is more suitable for homophily graphs, GPRGNN~\cite{chien2020adaptive} improves it by adaptively learning GPR weights. This method not only avoids over-smoothing but also enables adaptive filtering to handle heterophily graphs.
NCGNN~\cite{chen2024neighborhood} improves model capacity and training efficiency by pre-propagating node features, then employs a CNN-based approach to aggregate the propagated features, adaptively capturing contextual information to boost the performance.
% Our work is inspired by these decoupled designs, analyzing the influence relationships among labels on a graph from the perspectives of propagation (P) and transformation (T). 
Our work is inspired by above studies that decouple or disentangle~\cite{zhang2022model} the message passing process. 
We focus more on analyzing the influence correlations between labels during the P and T operations, rather than designing specific P and T operations to increase model performance~\cite{gasteiger2018predict,chien2020adaptive,ijcai2022p310} or enhance model capacity and scalability~\cite{chen2024neighborhood}.

% [1] Predict then propagate: Graph neural networks meet personalized pagerank, ICLR 2019
% gasteiger2018predict   
% appnp

% [4] Towards deeper graph neural networks, KDD 2020
% liu2020towards
% DAGNN

% GCNII不是 lightGCN SGC wu2019simplifying 没说，但是明显使用的

% [2] On the equivalence of decoupled graph convolution network and label propagation, WWW 2021
% dong2021equivalence
% PTA

% [3] Adaptive universal generalized pagerank graph neural network, ICLR 2021
% chien2020adaptive
% GPRGNN

% [5] Neighborhood Convolutional Graph Neural Network, KBS 2024
% chen2024neighborhood
% NCGNN

% We choose GCN~\citep{kipf2017semisupervised} as main backbone, ML-KNN~\citep{10.1016/j.patcog.2006.12.019} and latest work PLAIN~\citep{wang2023deep} as multi-label methods.
% , GAT~\citep{2017Graph}, GraphSAGE~\citep{hamilton2017inductive}
% \iclr{

% We analyze the influence between labels on graph separately within these two components and then aggregate them. 
% Note that even in some coupled GNNs, like GCN~\cite{kipf2017semisupervised} and GAT~\citep{2017Graph}, the operations of P and T inherently exist. Therefore, our analysis (in Sec.~\ref{sec:over}) of the influence between labels from these two step is also applicable to the coupled GNNs. 
% In short, this paper provides deeper insights into MLNC and introduces an effective MLNC model based on these analysis.
% }

